% [原文 p.321]
\chapter{S 类理论}

\ChapterAbstract{我们借助超弦理论与 M 理论的技术，按照 [366] 的开创性工作以及 [155, Sect. 3.2] 的详细处理，回顾一大类四维 $\mathcal{N}=2$ 理论的 Seiberg–Witten 曲线的构造。在简要概述超弦理论与 M 理论（着重其膜内容）之后，我们描述实现大量 $\mathcal{N}=2$ 规范理论族的 IIA 型膜构造。它们的 M 理论提升给出相应的 Seiberg–Witten 曲线，为系统地构造这些曲线提供了有力的框架。随后我们引入 class $\mathcal{S}$ 理论（S 类理论）的框架，取材于 [146, 339]。我们首先说明带四个味道的四维 $\mathcal{N}=2$ SU(2) SYM 的 $S$ 对偶性质如何产生广义 SU(2) 箭图，后者刻画一大类 $\mathcal{N}=2$ 超共形场论的弱耦合极限，即 $A_1$ 型的 S 类理论。这一构造自然地从 SU(2) 推广到 SU(3)，进而推广到 SU($N$)。}

本专著的核心对象——谱网络（将在下一章第 22 章从 S 类理论中 BPS 态计数的角度引入）——在高阶泰希米勒理论中已有重要应用，这已在第 III 部分讨论过。这一联系源于 S 类理论的一个界定性特征：它们与黎曼曲面有内在关联（S 代表“曲面”）。特别地，S 类理论的 Seiberg–Witten 曲线自然地作为另一条黎曼曲面的分歧覆盖（参看第 2 章）出现，后者称为 UV 曲线。此外，UV 曲线的泰希米勒空间（在第 6 章引入）及其映射类群（在第 4.5 节引入）都将获得具体的物理解释。

\section{超弦理论与 M 理论}

在弦理论中，基本对象是弦而非点状粒子。技术上，这意味着弦理论被微扰地定义为二维量子场论，其场是从一个黎曼曲面到（伪）黎曼靶空间 $X$（解释为时空）的映射。
% [原文 p.322]
在大于弦尺度的尺度上，弦看起来像普通粒子，其质量、电荷及其他性质由它的振动状态决定。这样，弦理论在低能下就在 $X$ 上定义了一个有效量子场论。

为了在 $X$ 中同时描述费米子与玻色子，需要考虑超对称弦理论。已知的超对称弦理论有五种：两种杂化弦理论（记为 HO 与 HE）、两种版本的 II 型弦理论（IIA 型与 IIB 型），以及 I 型超弦理论。下文中我们主要关注 II 型弦理论，尤其是 IIA 型及其强耦合极限，即 M 理论。

超弦理论的量子化对场施加了约束。解决这些约束的一个自然方式，是把靶空间 $X$ 的维数固定为所谓临界维数。II 型超弦理论的临界维数是 10。在低能下，IIA 型与 IIB 型超弦理论分别被 IIA 型与 IIB 型超引力很好地近似。类型 A（相应地，类型 B）在 10 维 Minkowski 时空 $\R^{1,9}$ 中具有 $\mathcal{N}=(1,1)$（相应地，$\mathcal{N}=(2,0)$）超对称。M 理论不是一种弦理论，而是一种膜（membrane）的理论；它自然地定义在 11 维，并在低能下被 11 维超引力很好地近似。

这五种弦理论由一张对偶之网相联系，也就是说，最好把它们看成同一个底层理论的不同描述，而不是彼此独立的若干理论。例如，T 对偶联系 IIA 型与 IIB 型弦理论。关于弦理论、超弦理论与 M 理论的教材很多，例如 [170, 171, 301, 302, 217, 43]。

IIA 型超引力的场包括一个标量场 $\phi$、一个度量 $g$、一个 2-形式 $B^{(2)}$、一个 1-形式 $\lambda^{(1)}$ 和一个 3-形式 $\lambda^{(3)}$，以及它们的超对称伙伴。这些场来源不同：$\phi$、$g$ 与 $B^{(2)}$ 来自闭弦的 Neveu–Schwarz（NS）扇区，而 $\lambda^{(1)}$ 与 $\lambda^{(3)}$ 来自 Ramond–Ramond（RR）扇区。场 $\lambda^{(1)}$ 是一个 U(1) 规范场，D0-膜与 D6-膜分别在其中带电力与磁荷。更高阶的规范场 $B^{(2)}$ 与基本弦（有时称为 NS1-膜）电耦合，与 NS5-膜磁耦合；而 $\lambda^{(3)}$ 与 D2-膜电耦合，与 D4-膜磁耦合。基本弦、NS5-膜、D2-膜与 D4-膜分别是 $X$ 中 2 维、6 维、3 维与 5 维的对象。所有这些对象都被认为存在于完整的 IIA 型超弦理论中，而不只是在其低能极限中。关键之处在于，膜是 BPS 的，意即它们保持 $X$ 中一部分超对称。

在超弦理论中，D-膜（其中 D 代表 Dirichlet）也可描述为开弦的边界条件；开弦在它们存在时的量子化，在足够低的能量下在膜上定义一个量子场论，称为膜上的世界体理论（worldvolume theory）。例如，$\R^{1,9}$ 中的 D4-膜支撑一个规范群为 U(1) 的 $5d$ $\mathcal{N}=2$ 超 Yang–Mills 理论。当 $N$ 张 D4-膜重合时，规范代数增强为 $\mathfrak{u}(N)$。另一方面，NS5-膜支撑的是自对偶反对称张量的规范理论，而不是
% [原文 p.323]
Yang–Mills 理论。对 $N$ 张重合的 NS5-膜，得到秩为 $N-1$ 的非阿贝尔自对偶反对称张量的超共形理论。

D-膜的界定性——基本弦可以终止于其上——经各种对偶序列的推广，意味着真正的膜常常可以终止于其他膜上 [337, 347]。例如，D4-膜可以终止于 NS5-膜上，从而在 $X$ 中构造保持 IIA 型超弦理论一部分超对称的 D4-膜与 NS5-膜之网。

在 11 维超引力中，场包括一个度量 $g$ 与一个 3-形式 $C^{(3)}$，以及它们的超伙伴。M 理论容许 BPS 的 M2-膜（在 $C^{(3)}$ 下带电荷）与 M5-膜（在 $C^{(3)}$ 下带磁荷）。M2-膜与 M5-膜可以终止于 M5-膜上。IIA 型超弦与 M 理论之间的关系如下：

\begin{proposition}[21.1] 弦耦合常数为 $\lambda$ 时，$X$ 上的 IIA 型理论等价于 $X\times S^1$ 上的 M 理论，其中 $S^1$ 的半径正比于 $\lambda$。\end{proposition}

圆周 $S^1$ 称为 M 理论圆。在极限 $\lambda\to\infty$ 下，M 理论圆退紧化。换言之，平坦 11 维 Minkowski 空间 $\R^{1,10}$ 中的 M 理论可理解为 IIA 型弦理论的无穷耦合极限。$X\times S^1$ 上的 M 理论中的膜下降为 $X$ 上的 IIA 型膜：基本弦对应于缠绕 $S^1$ 并沿 $X$ 中某曲面伸展的 M2-膜，而 $X$ 中的 D2-膜对应于不缠绕 $S^1$ 的 M2-膜。类似地，D4-膜对应于缠绕 $S^1$ 的 M5-膜，而 NS5-膜对应于不缠绕该圆周的 M5-膜。

相反，IIA 型超弦理论中的 D0-膜与 D6-膜似乎并不从任何 M 理论膜下降而来；事实上，D6-膜提升为 $X$ 上 $S^1$ 纤维化（M 理论即定义于其上）的奇点。存在 D6-膜时，IIA 型超弦理论中的膜构造提升为 M 理论膜的构造，但该构造不在 $X\times S^1$ 中，而在 $X$ 上一个带奇点的 $S^1$ 纤维化中。这类纤维化称为广义 Taub–NUT 几何。

\section{膜构造与 $\mathcal{N}=2$ 箭图}

\subsection{简单设定}

让我们先考虑 $\R^{1,9}$ 上 IIA 型超弦理论中 $N$ 张重合 D4-膜夹在两张 NS5-膜之间的构造，坐标为 $x^0,\dots,x^9$。具体地说，两张 NS5-膜沿方向 $0,1,2,3,4,5$ 伸展，（经典地）分别位于 $x^6=0$ 与 $x^6=L$；而 D4-膜沿 $0,1,2,3$ 以及沿 $x^6$ 的线段 $[0,L]$ 伸展，并位于 $x^{4,5}=0$。此外，所有膜都位于 $x^{7,8,9}=0$。该构造如图 21.1 左图所示。

\begin{fact}[21.2（见 {{{{[186, 366]}}}}）] 这样的膜构造保持 $\R^{1,9}$ 中 IIA 型超对称的四分之一，即在 32 个超荷中保持 8 个。\end{fact}

% [Tutor Note: IIA 型与 M 理论膜的词典可逐项核对：F1 = 缠绕 $S^1$ 的 M2，D2 = 不缠绕 $S^1$ 的 M2，D4 = 缠绕 $S^1$ 的 M5，NS5 = 不缠绕 $S^1$ 的 M5。维度核对：M2 世界体 3 维，缠绕 $S^1$ 后余 2 维 → F1 ✓；M5 世界体 6 维，缠绕 $S^1$ 后余 5 维 → D4 ✓。D0 对应 M 理论圆上的 Kaluza–Klein 动量模，D6 对应 $S^1$ 纤维化的奇点（Taub–NUT），二者都不出现在 $X\times S^1$ 的膜谱中。另 Fact 21.2 的“保持四分之一”即 32 个超荷中保 8 个（原超对称的 $1/4$）。]

% [原文 p.324]
\Fig[0.85\linewidth]{figures/fig_21_1.png}{21.1}{两个 IIA 型膜构造。}

由于 D4-膜沿 $x^6$ 的伸展有限，其世界体理论在足够低的能量下有效地约化为四维：具体地说，沿 $0,1,2,3$ 得到 $4d$ $\mathcal{N}=2$ $\mathfrak{u}(N)$（或者更确切地说，$\mathfrak{su}(N)$，我们稍后解释）超 Yang–Mills（SYM）理论。D4-膜上的世界体理论有效变为四维的能量阈值与 $L$ 成反比。

现在考虑一个相关构造：D4-膜沿 $x^{4,5}$ 方向以一般方式分开，如图 21.1 右图所示。D4-膜上的世界体理论成为规范群为 $\U(1)^N$（或者更确切地说，$\U(1)^{N-1}$）的 $4d$ $\mathcal{N}=2$ 理论，而不再是 $\U(N)$。该膜构造被解释为对应于 $4d$ $\mathcal{N}=2$ $\U(N)$ SYM 库仑分支上的一个特定点：分开 D4-膜堆叠相当于把 $\SU(N)$ 破缺到其 Cartan 子群 $\U(1)^{N-1}$。

\subsection{更一般的设定}

我们将关心已经引入的简单设定的推广，如图 21.2 所示。这些膜构造由位于固定 $x^6$ 值的 $k$ 张 NS5-膜以及夹在它们之间的 D4-膜组成。可以有半无限 D4-膜：一端附着在 NS5-膜上，另一端沿 $x^6$ 延伸到 $\pm\infty$。此外，还可以加入沿 $0,1,2,3,7,8,9$ 伸展的 D6-膜，而不破坏额外的超对称。虽然稍后会回到这一可能性，但我们先考虑只含 D4-膜与 NS5-膜的构造。

\subsection{反作用}

通过计入 NS5-膜内部 D4-端点带来的反作用，可以改进我们所考虑的膜构造的图示。这些端点是在 NS5-世界体中的带荷余维 2 对象。更确切地说，D4-端点是 NS5-膜上世界体理论的磁单极子或涡旋。

% [原文 p.325]
\Fig[0.75\linewidth]{figures/fig_21_2.png}{21.2}{一个更复杂的 IIA 型膜构造。}

这些源使 NS5-膜对数式地弯曲。这种弯曲的一个直接后果是：一般而言，NS5-膜并不位于固定的 $x^6$ 值。记 $v=x^4+\ii x^5\in\C$，考虑位于 $v=0$ 的一张 D4-膜终止于经典上位于 $x^6=0$ 的 NS5-膜。NS5-膜倾向于使其体积最小化。这一点，连同作为余维 2 源的 D4-端点的存在，意味着在离 D4-端点足够远的地方，NS5-膜沿 $x^6$ 的位置满足二维 Laplace 方程：
\[
\nabla^2x^6(v)=0.
\]
这意味着 NS5-膜沿 $x^6$ 的位置形如 $x^6(v)=k\ln(|v|)+c$，其中 $k$ 与 $c$ 为常数。

更一般地，对于位于 $v=a_1,\dots,a_{q_L}$ 并终止于 NS5-膜左侧的 $q_L$ 张 D4-膜，以及位于 $v=b_1,\dots,b_{q_R}$ 并终止于 NS5-膜右侧的 $q_R$ 张 D4-膜，NS5-膜沿 $x^6$ 的位置为：
\begin{equation}
x^6(v)=k\sum_{i=1}^{q_L}\ln|v-a_i|-k\sum_{j=1}^{q_R}\ln|v-b_j|+c.
\tag{21.1}
\end{equation}
因此，只有当 $q_L=q_R$ 时，NS5-膜才在 $|v|\to\infty$ 时趋于某个特定的 $x^6$ 值。只有在这类情形下才能给 NS5-膜指定一个确定的 $x^6$ 值，而这个值只表明 NS5-膜沿 $x^6$ 的位置在 $|v|\to\infty$ 时渐近地趋于该固定值。

% [Tutor Note: 式 (21.1) 是 IIA 型中的“经典”弯曲：$x^6$ 是实坐标，点源给出二维 Laplace 方程的解 $\ln|v-a|$。而 (21.5) 的 $s=(x^6+\ii x^{10})/R$ 是复坐标，用全纯的 $\ln(v-a)$；二者由 $\ln|v-a|=\operatorname{Re}\ln(v-a)$ 相联系，故 (21.5) 取实部即回到 (21.1) 的弯曲行为。核对：$|v|\to\infty$ 时 (21.1) 有有限极限当且仅当 $q_L=q_R$，这对应 (21.5) 中 $s$ 在无穷远处对数发散的抵消（单值性）。]

现在考虑 D4-膜沿 $v$ 的运动，即复参数 $a_i$ 与 $b_j$ 的变化。当 D4-膜移动时，它们诱导 NS5-膜的形变，后者又影响 D4-膜的动能。可以区分那些需要付出有限能量代价的 D4-膜堆叠形变（称为可归一形变）与其他形变。如 [366, Sect. 2] 所示，一个形变是可归一的，如果如下条件成立：
% [原文 p.326]
\begin{equation}
\sum_i a_i-\sum_j b_j=q\in\C,
\tag{21.2}
\end{equation}
其中 $q$ 为常数。该约束意味着跨过 NS5-膜的两组 D4-膜的质心之间没有相对位移。当只考虑可归一形变时（我们将如此），$q\in\C$ 充当构造中与该 NS5-膜相联系的常数。

\subsection{规范理论解释}

考虑如图 21.2 的一般构造，含 $k+1$ 张 NS5-膜，记为 NS5$_\alpha$（$\alpha=0,\dots,k$）；并且对每个 $1\le\alpha\le k$，有 $N_\alpha$ 张 D4-膜夹在相邻的 NS5$_{\alpha-1}$ 与 NS5$_\alpha$ 之间。我们把 $N_\alpha$ 的定义延拓为包含 $N_0$（构造左侧的半无限 D4-膜数目）与 $N_{k+1}$（构造右侧的半无限 D4-膜数目）。

粗略地说，对每个 $1\le\alpha\le k$，每堆 $N_\alpha$ 张 D4-膜对 D4-膜上世界体理论的规范群贡献一个 $\U(N_\alpha)$ 因子。延伸到无穷的 D4-膜不贡献规范群，但它们给出 $\U(N_1)$ 的反基本表示中的 $N_0$ 个超多重态，以及 $\U(N_k)$ 的基本表示中的 $N_{k+1}$ 个超多重态。更精确地说：

\begin{proposition}[21.3] 式 (21.2) 表明，对每个 $2\le\alpha\le n-1$，乘积 $\U(N_{\alpha-1})\times\U(N_\alpha)$ 中的对角 $\U(1)$ 因子被冻结。因此，D4-膜上世界体理论的规范群 $G$ 为：
\[
G=\prod_{\alpha=1}^{k}\SU(N_\alpha).
\]
\end{proposition}

% [Tutor Note: 原文此处写 “for each $2\le\alpha\le n-1$”，但本节记号只有 $k+1$ 张 NS5-膜以及 $N_0,\dots,N_{k+1}$，并未出现 $n$；按语境（每条相邻因子对 $\U(N_{\alpha-1})\times\U(N_\alpha)$ 的对角 $\U(1)$ 被冻结）应为 $2\le\alpha\le k$，共冻结 $k-1$ 个对角 $\U(1)$，从而留下 $\prod_{\alpha=1}^{k}\SU(N_\alpha)$。照译原文并在此指出。]

半无限 D4-膜仍然给出 $\SU(N_1)$ 的反基本表示中的 $N_0$ 个超多重态，以及 $\SU(N_k)$ 的基本表示中的 $N_{k+1}$ 个超多重态。这些超多重态分别贡献 $\U(N_0)$ 与 $\U(N_{k+1})$ 味道因子。例如，$\SU(N)$ 基本表示中的 $M$ 个超多重态可记为：
\[
(Q^a_i,\tilde Q^a_i),\qquad a=1,\dots,N,\ \text{以及}\ i=1,\dots,M.
\]
这里的味道对称性\footnote{我们隐含地假设 $N$ 大于 2；对 $N=2$，相应的味道因子会增强为 $\SO(2N)$（参看命题 20.19）。}是 $\U(M)\simeq\SU(M)\times\U(1)$，其中单因子 $\SU(M)$ 作用在味道指标 $i$ 上。具体地说，$Q$ 在 $\SU(M)$ 的反基本表示中变换，而 $\tilde Q$ 在基本表示中变换。在阿贝尔因子 $\U(1)$ 下，$Q$ 的荷为 $+1$，$\tilde Q$ 的荷为 $-1$。

% [原文 p.327]
类似地，对每张 NS5-膜 NS5$_\alpha$（$2\le\alpha\le k-1$），存在一个处于 $\SU(N_{\alpha-1})\times\SU(N_\alpha)$ 的双基本表示中的超多重态，即：
\[
(Q^a_b,\tilde Q^b_a),\qquad a=1,\dots,N_\alpha,\ \text{以及}\ b=1,\dots,N_{\alpha-1}.
\]
这类双基本超多重态可以是有质量（BPS）超多重态，其裸质量等同于式 (21.2) 中出现的常数 $q_\alpha$。把这类规范理论表示为 $\mathcal{N}=2$ 箭图是很方便的，如图 21.3 所示。

\Fig[0.8\linewidth]{figures/fig_21_3.png}{21.3}{把所论理论表示为一个线性 $\mathcal{N}=2$ 箭图。}

\begin{definition}[21.4] 一个 $\mathcal{N}=2$ 箭图是有限连通图，含两类节点：圆节点对应 SU 规范因子，方节点代表 U 味道因子。每个节点（无论圆或方）都标以一个正整数：标以 $N$ 的圆节点描绘一个 $\SU(N)$ 规范因子，而标以 $N$ 的方节点描绘一个 $\U(N)$ 味道因子。边无向，或连接两个圆节点，或连接一个圆节点与一个方节点；它们代表双基本超多重态。\end{definition}

更一般地，规范群为乘积 $G=\prod G_i$、且带规范因子双基本表示中的手征超场的 $\mathcal{N}=1$ 规范理论，可以表示为如定义 5.1 所定义的箭图，即有限有向图。在这一框架下，箭图的节点代表因子 $G_i$，而从 $G_i$ 到 $G_j$ 的箭头表示一个双基本手征超场。$G_i$ 的伴随表示中的手征超场由节点 $G_i$ 处的一个圈表示。可以修改定义以把味道群也纳入。

表示 $R$ 中的一个超多重态对应于一对手征超场，一个在 $R$ 中，另一个在共轭表示 $\bar R$ 中。因此，规范群为 $G=\prod G_i$ 且含双基本超多重态的 $\mathcal{N}=2$ 规范理论的 $\mathcal{N}=1$ 箭图，其箭头总是成对出现且方向相反。按照惯例，在 $\mathcal{N}=2$ 箭图中，这些对子被画成一条无向边。每个 $\mathcal{N}=2$ 矢量超场也包含一个伴随表示中的 $\mathcal{N}=1$ 手征超场，后者在 $\mathcal{N}=2$ 箭图图中通常省略。

在某些箭图的定义中，人们排除圈以及方向相反的一对箭头。然而，圈与这样的对子在 $\mathcal{N}=2$ 箭图中很常见。这与带势箭图（quiver with potential）的结构相容，后者允许这些元素（例如见 [105]）。

从一般原理出发，可以期望：对应于夹在经典上位于 $x^6_{\alpha-1}$ 与 $x^6_\alpha$ 的 NS5-膜之间的 $N_\alpha$ 张 D4-膜堆叠的 $\SU(N_\alpha)$ 规范理论的规范耦合，由下式给出：

% [原文 p.328]
\begin{equation}
\frac{1}{g_\alpha^2}=\frac{x^6_\alpha-x^6_{\alpha-1}}{\lambda},
\tag{21.3}
\end{equation}
其中 $\lambda$ 是 IIA 型弦耦合。然而我们已看到，一般而言 $x^6_{\alpha-1}$ 与 $x^6_\alpha$ 依赖于复坐标 $v$，后者参数化 D4-膜在 $v=x^4+\ii x^5$ 平面中的位置。如前，设 $N_{\alpha-1}$ 为位于 $x^6_{\alpha-1}$ 的 NS5-膜左侧的 D4-膜数目，$N_{\alpha+1}$ 为位于 $x^6_\alpha$ 的 NS5-膜右侧的 D4-膜数目。若 $N_{\alpha-1}$ 与 $N_{\alpha+1}$ 都小于 $N_\alpha$，则在 $v$ 很大时，$x^6_\alpha-x^6_{\alpha-1}$ 按 $(2N_\alpha-N_{\alpha-1}-N_{\alpha+1})k\ln|v|$ 对数式增长。

从 $\SU(N_\alpha)$ 规范理论的角度看，相邻的 $N_{\alpha-1}$ 与 $N_{\alpha+1}$ 张 D4-膜堆叠对规范耦合 $g_\alpha$ 的重整化所作的贡献，与 $N_{\alpha-1}+N_{\alpha+1}$ 个基本超多重态相同。因此：

\begin{proposition}[21.5] 修正后的方程
\begin{equation}
\frac{1}{g_\alpha^2}(v)=\frac{x^6_\alpha(v)-x^6_{\alpha-1}(v)}{\lambda}
\tag{21.4}
\end{equation}
反映了规范群 $\SU(N_\alpha)$ 的单圈 beta 函数，其中 $b_1=2N_\alpha-N_{\alpha-1}-N_{\alpha+1}$。相应地，距离 $|v|$ 被解释为能标：$|v|$ 大对应于 UV，此时 $g_\alpha^2$ 趋于 0；而 $|v|$ 小对应于 IR。\end{proposition}

\subsection{更一般的 $\mathcal{N}=2$ 箭图规范理论}

可以加入沿 $x^{0,1,2,3,7,8,9}$ 方向伸展、位于固定 $x^6$ 值与 $v=x^4+\ii x^5$ 值的无限 D6-膜，而不比只含 D4-膜与 NS5-膜的构造破坏更多超对称。因此，D4-膜上的世界体理论仍是 $4d$ $\mathcal{N}=2$ 理论。加入 D6-膜会为某些 $\SU(N_\alpha)$ 规范因子引入基本超多重态，导致如图 21.4 所示的更一般的线性箭图规范理论。图 21.4 的箭图对应于在 NS5-膜 NS5$_{\alpha-1}$ 与 NS5$_\alpha$ 之间放置 $F_\alpha$ 张 D6-膜的膜构造。

\Fig[0.8\linewidth]{figures/fig_21_4.png}{21.4}{更一般的线性 $\mathcal{N}=2$ 箭图。}

也可以考虑把 $x^6$ 方向取为紧致的构造，即构成一个圆周。这样的构造中不能有半无限 D4-膜。
% [原文 p.329]
编码 D4-膜上低能理论结构的箭图构成一个闭合环路，称为项链箭图。项链箭图是非线性的 $\mathcal{N}=2$ 箭图之例：规范群是 $k$ 个 $\SU(N_\alpha)$ 因子的乘积\footnote{D4-膜上世界体理论的规范群还含一个 $\U(1)$ 因子，即 $\prod\U(N_\alpha)$ 的对角 $\U(1)$ 子群。该 $\U(1)$ 因子是自由的，并在 IR 中退耦。}，每个因子对应夹在相邻 NS5-膜之间的一堆 D4-膜。此外，若在此构造中引入 D6-膜，它们如前贡献额外的味道群。所得理论称为椭圆模型，理由稍后就会清楚。

\noindent\textbf{共形箭图}

在拉氏量 $4d$ $\mathcal{N}=2$ 理论中，每个规范耦合的重整化在单圈是微扰精确的。图 21.4 的线性箭图理论中规范因子 $\SU(N_\alpha)$ 的单圈 beta 函数系数由下式给出：
\[
b_{1,\alpha}=2N_\alpha-N_{\alpha-1}-N_{\alpha+1}-F_\alpha.
\]
若对每个 $\alpha$ 都有 $N_\alpha\ge N_{\alpha-1}+N_{\alpha+1}$，则可以调整 $F_\alpha$ 使得 $b_{1,\alpha}=0$。满足该条件的箭图被认为是共形的，因此称为共形箭图。

% [Tutor Note: 可用标准公式核对：$\SU(N)$ 带 $N_f$ 个基本味道的单圈系数为 $b_1=2N-N_f$。这里 $\SU(N_\alpha)$ 的有效味道数应为 $N_{\alpha-1}+N_{\alpha+1}+F_\alpha$（相邻两堆各对“味道”有贡献，D6-膜贡献 $F_\alpha$ 个），故 $b_{1,\alpha}=2N_\alpha-(N_{\alpha-1}+N_{\alpha+1})-F_\alpha$，与原文一致。共形条件 $b_{1,\alpha}=0$ 即 $F_\alpha=2N_\alpha-N_{\alpha-1}-N_{\alpha+1}$（需要 $N_\alpha\ge N_{\alpha-1}+N_{\alpha+1}$ 才能取 $F_\alpha\ge0$），正是正文所述；例 21.6 中 $N_\alpha=N$、$F_\alpha=0$ 给出 $b_{1,\alpha}=0$ ✓。]

\begin{example}[21.6] 取 $N_\alpha=N$（$0\le\alpha\le k+1$）以及 $F_\alpha=0$（$1\le\alpha\le k$），就定义了规范群为 $\SU(N)^k$ 的共形箭图，其中第一个与最后一个规范群各有 $N$ 个基本超多重态。当 $k=1$ 时，该理论是带 $2N$ 个基本味道的 $\SU(N)$ SYM 理论。\end{example}

\Fig[0.95\linewidth]{figures/fig_21_5.png}{21.5}{平衡箭图。}

% [原文 p.330]
\begin{example}[21.7] 图 21.5 所示的箭图——它们对应于仿射 $A$、$D$、$E$ Dynkin 图\footnote{事实上，这些 ADE 箭图是唯一不带味道群的共形箭图 [158]。这来自如下观察：共形性条件定义的是仿射 Cartan 矩阵。}——是共形的。它们有时称为平衡箭图。$\widehat A_{k-1}$ 的情形是项链箭图。$\widehat D_n$ 型的仿射箭图也可以通过引入 NS-orientifold 平面在 IIA 型超弦中实现 [224]。某些 $4d$ $\mathcal{N}=2$ 箭图尚无已知的膜实现；例如，目前对 $\widehat E_6$、$\widehat E_7$ 或 $\widehat E_8$ 箭图理论都不知道任何膜构造。\end{example}

\section{M 理论提升与 Seiberg–Witten 曲线}

我们已经识别出一大类 $4d$ $\mathcal{N}=2$ 规范理论，它们可实现为 IIA 型超弦中 D4-膜上的世界体理论，其构造涉及 D4-膜、NS5-膜以及可能的 D6-膜。当每一堆“规范” D4-膜由重合的膜组成时，理论是一个规范群为若干单因子之积的 $\mathcal{N}=2$ 箭图。相邻堆沿 $v$ 的位置差编码双基本超多重态的裸质量，而半无限 D4-膜或 D6-膜的分离编码（反）基本超多重态的裸质量。这类构造的可归一形变（D4-膜在 $v=x^4+\ii x^5$ 平面中分开）对应于这些 $\mathcal{N}=2$ 箭图的库仑分支。下面我们将解释，$C$ 上的低能动力学（编码在 $C$ 上的 Seiberg–Witten 纤维化中）如何从弦理论、更具体地说从 IIA 型膜构造到 M 理论的提升中自然地出现。

\subsection{M 理论提升}

IIA 型超弦理论的 M 理论极限对应于 $\lambda\to\infty$。式 (21.3) 与 (21.4) 表明，如果同时重标度 $x^6$ 方向，就可以在不改变 $4d$ 规范耦合的情形下取此极限。因此，能实现为 IIA 型超弦中膜构造的 $\mathcal{N}=2$ 箭图，同样能实现为 M 理论中的膜构造，称为 IIA 型构造的 M 理论提升。

记 $x^{10}\sim x^{10}+2\pi R$ 为 M 理论圆 $S^1$ 上的周期坐标（其中 $R$ 正比于 $\lambda$）。D4-膜与 NS5-膜都从 M 理论中的 M5-膜下降而来：沿 $x^{0,1,2,3,6}$ 伸展的 D4-膜从沿 $x^{0,1,2,3,6,10}$ 伸展的 M5-膜下降；而沿 $x^{0,1,2,3,4,5}$ 伸展、置于特定 $x^6$ 值的 NS5-膜，从沿同样坐标伸展、置于固定 $x^6$ 与 $x^{10}$ 值的 M5-膜下降。例如，与两张 NS5-膜相交的 $N$ 张平行 D4-膜提升为沿 $x^{0,1,2,3,6,10}$ 伸展的 $N$ 张 M5-膜（因而缠绕由 $x^{6,10}$ 参数化的柱面），以及两张沿 $x^{0,1,2,3,4,5}$ 伸展、置于柱面上特定点的 M5-膜，如图 21.6 所示。该 M 理论构造对应于库仑分支原点处带 $F=2N$ 个基本味道的 $\mathcal{N}=2$ $\SU(N)$ SYM。

% [原文 p.331]
\Fig[0.9\linewidth]{figures/fig_21_6.png}{21.6}{把 IIA 型膜构造提升到 M 理论。}

在 D4-膜夹在 NS5-膜之间的一般 IIA 型构造中，NS5-膜沿 $x^6$ 方向的位置依赖于 $v$。在 M 理论中对应的陈述是：任何 NS5-膜的 M5-膜提升沿 $x^6+\ii x^{10}$ 的位置也依赖于 $v$。为使 M 理论构造具有超对称，$x^6+\ii x^{10}$ 必须全纯地依赖于 $v$。这导致式 (21.1) 的如下推广，其中 $s:=(x^6+\ii x^{10})/R$：
\begin{equation}
s(v)=\sum_{i=1}^{q_L}\ln(v-a_i)-\sum_{j=1}^{q_R}\ln(v-b_j)+c.
\tag{21.5}
\end{equation}
令 $\tau_\alpha=4\pi\ii/g_\alpha^2+\theta_\alpha/(2\pi)$ 为 $\SU(N_\alpha)$ 规范理论（在夹于两张 NS5-膜 NS5$_{\alpha-1}$ 与 NS5$_\alpha$ 之间的 $N_\alpha$ 张 D4-膜世界体上）通常的复化规范耦合，则式 (21.4) 中的表达式推广为：
\[
\ii\tau_\alpha(v)=s_\alpha(v)-s_{\alpha-1}(v).
\]
这一推广被解释为耦合 $\tau_\alpha$ 的微扰重整化。可以看出，向 M 理论的提升自然地把 theta 角纳入第 21.2 节的讨论之中。

如图 21.2 所示的一般 D4-膜与 NS5-膜构造的 M 理论提升，必须对应于复二维空间 $Q=\C_v\times\C^\times_s$ 中的一条黎曼曲面 $\Sigma$，才能保持与原 IIA 型膜构造相同数量的超对称。IIA 型超弦理论中的可归一形变提升为 M 理论中的可归一形变 [155]。在 M 理论圆收缩的 IIA 型极限下，$\Sigma$ 必须还原出原来的膜构造。两个关键洞察如下：
\begin{itemize}
  \item IIA 型膜构造中出现的奇点在 $\Sigma$ 中一般被磨光。这意味着关于 M 理论低能有效
% [原文 p.332]
近似的知识就足以确定 D4-膜上世界体理论的低能物理。
  \item 黎曼曲面 $\Sigma$ 正是 D4-膜上 $4d$ $\mathcal{N}=2$ 理论的 Seiberg–Witten 曲线。
\end{itemize}

第二点可论证如下。M5-膜上的世界体理论描述一个自对偶 2-形式规范场 $\beta$，即满足 $T=\mathrm{d}\beta=\star_6\mathrm{d}\beta$ 的 2-形式。设 M5-膜紧化在黎曼曲面 $\Sigma$ 上，从而其形式为 $\R^{1,3}\times\Sigma$。在所关心的情形中，$\Sigma$ 可通过加入有限多个点自然地紧化，得到亏格为 $g$ 的紧黎曼曲面 $\bar\Sigma$。于是 $\beta$ 在 $\bar\Sigma$ 上的零模有效地给出 $\R^4$ 上的 $g$ 个阿贝尔规范场，其复化耦合常数由 Jacobian $J(\bar\Sigma)$ 描述。这来自如下事实：$\bar\Sigma$ 上每个调和 1-形式 $\Lambda$ 都提供一种方式，把 $\R^{1,3}$ 中 Maxwell 方程的解嵌入为 $\R^{1,3}\times\bar\Sigma$ 中 $T$ 的解 [351]。

换言之，D4-膜世界体理论库仑分支上的低能物理由黎曼曲面 $\Sigma$ 及其 Jacobian 簇 $J(\bar\Sigma)$ 描述。$\Sigma$ 的复结构由具体的膜构造决定，等价地，由该构造所对应的库仑分支上的点决定。Jacobian 簇 $J(\bar\Sigma)$ 编码所得 $\R^{1,3}$ 中阿贝尔规范场的低能有效作用量。对子 $(\Sigma,J(\bar\Sigma))$ 代表 Hitchin 可积系统在库仑分支上特定点处的纤维，如第 20.3 节所讨论。这把 $\Sigma$ 等同于所论理论在其库仑分支上对应于给定膜构造的那个特定点处的 Seiberg–Witten 曲线。

\begin{remark}[21.8] 线性 $\mathcal{N}=2$ 箭图的 Seiberg–Witten 微分 $\lambda$ 可以等同于 $Q=\C_v\times\C^\times_s$ 上的 Liouville 1-形式（把 $Q$ 视为 $\C^\times_s$ 的全纯余切丛）。用坐标写出：
\[
\lambda=v\frac{\mathrm{d}t}{t},
\]
其中 $t=\exp(-s)$ 是 $\C^\times_s$ 上一个整体定义的复坐标。\end{remark}

\subsection{Seiberg–Witten 曲线的计算}

考虑带两张 NS5-膜（记为 NS5$_0$ 与 NS5$_1$）的 IIA 型膜构造：NS5$_0$ 左侧有 $N_0$ 张半无限 D4-膜，NS5$_0$ 与 NS5$_1$ 之间有 $N$ 张 D4-膜，NS5$_1$ 右侧有 $N_2$ 张半无限 D4-膜。该构造如图 21.7 所示，其中我们示意了 NS5-膜在 $v$ 很大时的对数弯曲（在 $v$ 与 $x^6$ 很小时，IIA 型膜描述不准确：目前没有关于 NS5-膜内部 D4-端点的精确描述）。

由于 $s=(x^6+\ii x^{10})/R$ 是多值函数，我们用 $t=\ee^{-s}$ 而非 $s$ 来参数化 $\C^\times_{6,10}$。图 21.7 中构造的 M 理论提升由一张沿 $x^{0,1,2,3}$ 伸展、缠绕 $\C_v\times\C^\times_t$ 中黎曼曲面 $\Sigma$ 的 M5-膜组成。

% [原文 p.333]
\Fig[0.6\linewidth]{figures/fig_21_7.png}{21.7}{IIA 型膜构造。}

这条曲线 $\Sigma$——它是 D4-膜上世界体理论（此情形即带 $N_0+N_2$ 个味道的 $\SU(N)$ SYM）的 Seiberg–Witten 曲线——由一个多项式方程描述：
\[
F(t,v)=0.
\]
$F$ 的确切形式可如下导出：
\begin{itemize}
  \item 在 $v$ 的一般固定值 $v_0$ 处，$F(t,v_0)$（视为 $t$ 的多项式）的根对应于 IIA 型构造中 NS5-膜的位置。因此，对图 21.7 所示的构造，$F(t,v)$ 必须是 $t$ 的二次多项式，即：
\[
F(t,v)=A(v)t^2+B(v)t+C(v).
\]
  \item 在 $t$ 的固定值 $t_0$ 处（使得 $x^6=-R\ln|t|$ 落在 NS5$_0$ 与 NS5$_1$ 沿 $x^6$ 方向的经典位置之间），我们预期 $F(t_0,v)$ 有 $N$ 个根，对应于夹在 NS5$_0$ 与 NS5$_1$ 之间的 D4-膜的位置。因此，$A(v)$、$B(v)$ 与 $C(v)$ 是次数至多为 $N$ 的多项式。
  \item 在 $C(v)$ 的一个零点 $v_0$ 处，$F(t,v_0)$ 的一个根趋于 $t=0$，即 $x^6=\infty$；这可解释为在 $v=v_0$ 处存在一张位于 NS5$_1$ 右侧的半无限 D4-膜。因此 $C(v)$ 的次数为 $N_2$。
  \item 在 $A(v)$ 的一个零点 $v_0$ 处，$F(t,v_0)$ 的一个根趋于 $t=\infty$，即 $x^6=-\infty$；这可解释为在 $v=v_0$ 处存在一张位于 NS5$_0$ 左侧的半无限 D4-膜。因此 $A(v)$ 的次数为 $N_0$。
\end{itemize}

\begin{example}[21.9] 对纯 $4d$ $\mathcal{N}=2$ $\mathfrak{su}(N)$ SYM，可重标度 $t$ 使常数多项式 $A$ 与 $C$ 都等于 1。这给出 Seiberg–Witten 曲线：
\begin{equation}
t^2+B(v)t+1=0,
\tag{21.6}
\end{equation}
其中 $B(v)$ 是次数为 $N$ 的多项式。特别地，对 $N=2$，令 $z=t$ 并把 $v$ 重标度/平移为新坐标 $x$，就得到如下形式的等价方程：
% [原文 p.334]
\[
\Lambda^2z+\frac{\Lambda^2}{z}=x^2-u,
\]
这正是式 (20.7) 所给的纯 $\mathfrak{su}(2)$ SYM 的 Seiberg–Witten 曲线。\end{example}

现在令 $\tilde t=t+B(v)/2$，把方程改写为：
\begin{equation}
\tilde t^2=\frac{B(v)^2}{4}-1.
\tag{21.7}
\end{equation}
通过适当地重标度与平移 $v$，可进一步假设 $B(v)$ 具有如下形式：
\[
B(v)=v^N+u_2v^{N-2}+u_3v^{N-3}+\cdots+u_N.
\]
以该 $B(v)$ 代入的式 (21.6) 是纯 $\mathfrak{su}(N)$ SYM 的 Seiberg–Witten 曲线的标准表达。

\begin{remark}[21.10] 注意当 $|t|$ 很大时，式 (21.6) 的根近似由 $t\propto v^N$ 给出；类似地，当 $|t|$ 很小时，式 (21.6) 的根近似由 $t\propto v^{-N}$ 给出。在固定 $t$ 时，$F(t,v)$（视为 $v$ 的多项式）对任何一般的 $t$ 都有 $N$ 个根，而不只是对 NS5-膜沿 $x^6$ 的经典位置之间的那些 $t$：$t\propto v^{\pm N}$ 与式 (21.5) 的渐近形式相符，后者反映了 NS5-膜的对数弯曲。\end{remark}

\begin{example}[21.11] $N_0=0$、$N=2$、$N_2=1$ 的情形对应于带一个基本超多重态的 $\SU(2)$ SYM 理论。其 Seiberg–Witten 曲线写作：
\[
t+B(v)+\frac{\alpha v+\beta}{t}=0,
\]
其中 $B(v)$ 是次数为 2 的多项式，$\alpha$ 与 $\beta$ 为常数。通过适当地重标度与平移 $v$ 与 $t$，该方程可写成式 (20.13) 的形式：
\[
\Lambda_1^2z+\frac{2\Lambda_1(x-m_1)}{z}=x^2-u.
\]
\end{example}

\begin{example}[21.12] 更一般地，若 $N_0=0$、$N$ 任意而 $N_2=F$，则式 (21.7) 变为：
\[
\tilde t^2=\frac{B(v)^2}{4}-f\prod_{j=1}^{F}(v-m_j),
\]
其中 $f$ 是复常数，并且（可能在平移 $v$ 之后）：
\[
B(v)=e(v^n+u_2v^{n-2}+\cdots+u_{n-1}v+u_n),
\]
其中 $e$ 是另一个复常数。

对 $F\ne 2N$，可以重标度 $\tilde t$ 与 $v$ 以令 $e=f=1$。相反，对 $F=2N$，只能固定 $e$ 与 $f$ 的一个组合，留下一个未定的自由度。这个额外参数对应于理论的耦合常数，它在
% [原文 p.335]
$F=2N$ 时恰是边缘的。在膜图像中，两张 NS5-膜在 $|v|\to\infty$ 时变为平行。$|v|$ 很大时 NS5-膜之间的距离趋于一个常数，被解释为理论的无量纲 UV 规范耦合。\end{example}

\begin{remark}[21.13] 当 $N_0+N_2>2N$ 时，NS5-膜在 $|v|\to\infty$ 时也保持平行。这可解释为：相应 IIA 型膜构造的 M 理论提升定义了这些红外自由理论的 UV 完备化。\end{remark}

涉及 D4-膜与 NS5-膜的更一般构造可用同样的方法研究，从而为一大类线性 $4d$ $\mathcal{N}=2$ 箭图理论导出 Seiberg–Witten 曲线。这些曲线取为 $Q=\C_v\times\C^\times_t$ 中的黎曼曲面，而 Seiberg–Witten 微分由 Liouville 1-形式表示：
\[
\lambda=x\frac{\mathrm{d}t}{t}.
\]
最一般的线性箭图通过把 D6-膜纳入这些构造而得到。含 D6-膜的构造的 M 理论提升导致更复杂的几何，具体地说，复空间 $Q$ 被一个 Taub–NUT 几何取代。Seiberg–Witten 曲线仍是 $Q$ 中的一条黎曼曲面。也可以计算项链箭图的 Seiberg–Witten 曲线。在此情形，$Q$ 对应于椭圆曲线 $E$ 的全纯余切丛，故有椭圆模型之名。

\begin{example}[21.14] 带伴随表示中无质量超多重态的四维 $\mathcal{N}=2$ $\mathfrak{su}(N)$ SYM 理论——即 $4d$ $\mathcal{N}=4$ $\mathfrak{su}(N)$ SYM——由沿 $x^{0,1,2,3,6}$ 方向伸展的 $N$ 张 D4-膜描述，其中 $x^6$ 方向紧化在长度为 $2\pi L$ 的圆周上。一张 NS5-膜位于该圆周上的某点。当把此构造提升到 M 理论时，$x^6$ 与 $x^{10}$ 方向结合成一条椭圆曲线 $E$，其定义使得当 $x^6$ 平移 $2\pi L$ 时，$x^{10}$ 坐标也平移 $\theta R$（对某个 $\theta\in[0,2\pi[$）。此构造中的 IIA 型膜变成一张 M5-膜，它缠绕空间 $E\times\C_v\simeq T^*_{1,0}E$ 中 $E$ 的一个 $N$ 次分歧覆盖 $\Sigma$。曲面 $\Sigma$ 对应于 $4d$ $\mathcal{N}=4$ $\mathfrak{su}(N)$ SYM 的 Seiberg–Witten 曲线。

为给伴随超多重态引入裸质量 $m$，必须考虑 $\C_v$ 在圆周 $S^1_6$ 上的非平凡仿射纤维化，使得当 $x^6$ 平移 $2\pi L$ 时，坐标 $v$ 变换为 $v\mapsto v+m$（其中 $m\in\C$）。把这一修改后的膜构造提升到 M 理论，就得到 $4d$ $\mathcal{N}=2^*$ $\mathfrak{su}(N)$ 理论的 Seiberg–Witten 曲线，与 [114] 的结果一致。\end{example}

\subsection{BPS 态}

在 Seiberg–Witten 曲线的这一几何实现中，BPS 态与沿 $x^0$（$\R^{1,3}_{0,1,2,3}$ 中 BPS 粒子的世界线）伸展、并缠绕 $Q$ 中一条边界落在 Seiberg–Witten 曲线 $\Sigma$ 上的黎曼曲面 $D$ 的 M2-膜相联系。若 $D$ 在其相对于边界的同调类中使体积最小，则其量子化给出一个 BPS 态 [265]。我们将在第 22 章发展这一几何视角。

% [原文 p.336]
\section{带 $F=4$ 个味道的 SU(2) SYM 再探}

\subsection{电磁对偶}

回忆第 20.2 节关于带 $F=4$ 个无质量基本味道的 $4d$ $\mathcal{N}=2$ SU(2) SYM 的 $S$ 对偶的讨论。对偶群 $\SL_2(\Z)$ 联合作用在库仑分支上的特殊坐标 $a$ 与 $a_D$ 上，在此特殊情形它们满足（参看注 20.25）：
\[
\begin{aligned}
a(u)&=\sqrt{u},\\
a_D(u)&=2\tilde\tau_{\mathrm{cl}}a(u),
\end{aligned}
\]
并按三重性（triality）置换味道群 $\SO(8)$（更确切地说，$\Spin(8)$）的表示。

让我们聚焦于味道群的子群 $\Spin(4)\times\Spin(4)$，它产生于把四个基本超多重态分成两对。利用同构 $\mathfrak{so}(4)\simeq\mathfrak{su}(2)\times\mathfrak{su}(2)$，把第一个 $\Spin(4)$ 因子改写为 $\SU(2)_a\times\SU(2)_b$，把第二个改写为 $\SU(2)_c\times\SU(2)_d$。记 $\mathbf{2}_a,\mathbf{2}_b,\mathbf{2}_c,\mathbf{2}_d$ 为每个 $\SU(2)$ 副本的基本表示。

\begin{proposition}[21.15] 关于子群 $\SU(2)_a\times\SU(2)_b\times\SU(2)_c\times\SU(2)_d$，$\Spin(8)$ 的三个八维表示分解为：
\begin{equation}
\begin{aligned}
\mathbf{8}_v&\simeq(\mathbf{2}_a\otimes\mathbf{2}_b)\oplus(\mathbf{2}_c\otimes\mathbf{2}_d),\\
\mathbf{8}_s&\simeq(\mathbf{2}_a\otimes\mathbf{2}_c)\oplus(\mathbf{2}_b\otimes\mathbf{2}_d),\\
\mathbf{8}_c&\simeq(\mathbf{2}_a\otimes\mathbf{2}_d)\oplus(\mathbf{2}_b\otimes\mathbf{2}_c).
\end{aligned}
\tag{21.8}
\end{equation}
\end{proposition}

这里 $\mathbf{8}_v$、$\mathbf{8}_s$ 与 $\mathbf{8}_c$ 分别表示 $\Spin(8)$ 的矢量、旋量与共轭旋量表示。$\SL_2(\Z)$ 的元素按它们模 2 的约化置换 $\mathbf{8}_v$、$\mathbf{8}_s$ 与 $\mathbf{8}_c$——回忆 $\SL_2(\Z_2)\simeq\mathfrak{S}_3$。因此，对偶群 $\SL_2(\Z_2)$ 以一切可能的方式置换 $\SU(2)$ 因子。

对偶群对 $\SU(2)$ 味道因子的作用，由图 21.8 中精细化的箭图记号显示。如前，每个标以 2 的圆节点代表一个 $\SU(2)$ 规范因子，而每个标以 2 的方节点表示一个 $\SU(2)$ 味道因子\footnote{与上文相反，这里的方节点对应 SU 而非 U 的味道因子。}。三价顶点对应于若干对基本超多重态：例如，图 21.8 中最左边的三价顶点代表 $\SU(2)$ 规范群基本表示中的两个超多重态。与这些超多重态相联系的味道对称性是 $\SU(2)_a\times\SU(2)_b$。

% [原文 p.337]
\Fig[0.95\linewidth]{figures/fig_21_8.png}{21.8}{带 4 个基本味道的 $4d$ $\mathcal{N}=2$ SU(2) SYM 的 $S$ 对偶。}

图 21.8 中的箭图代表带 $F=4$ 个基本超多重态的四维 $\mathcal{N}=2$ SU(2) SYM 理论的弱耦合极限。最初，理论由图 21.8 最左边的箭图描绘，此时基本超多重态在 $\Spin(8)$ 的矢量表示 $\mathbf{8}_v$ 中变换。增大 $\SU(2)$ 规范群的规范耦合——在极限 $\tau\to0$ 下——并施加变换 $S$，得到对该理论更好的描述，即图 21.8 中间的箭图。在这一对偶描述中，规范耦合由 $\tau_D=-1/\tau$ 给出，意味着新的规范群是弱耦合的。

过渡到中间箭图之所以发生，是因为 $S$ 交换 $\Spin(8)$ 的 $\mathbf{8}_v$ 与 $\mathbf{8}_s$ 表示，同时保持 $\mathbf{8}_c$ 表示不变。这由 $S$ 模 2 的约化推出\footnote{如注 20.25 下方所讨论，$\Spin(8)$ 的表示由它们在中心 $\mathcal{Z}(\Spin(8))=\Z_2\times\Z_2$ 下的荷分类。具体地说，$\mathbf{8}_v$ 对应 $(1,0)$，$\mathbf{8}_s$ 对应 $(0,1)$，$\mathbf{8}_c$ 对应 $(1,1)$。}。

\noindent\textbf{有质量超多重态}

编码基本超多重态质量的矩阵 $m$ 属于 $\mathfrak{so}_8(\C)$ 的 Cartan 子代数。关于子群 $\SU(2)_a\times\SU(2)_b\times\SU(2)_c\times\SU(2)_d\subset\Spin(8)$，矩阵 $m$ 分解为四个复质量参数 $m_a$、$m_b$、$m_c$、$m_d$，各与其中一个 $\SU(2)$ 因子相联系。若基本超多重态的原质量记为 $m_1,m_2,m_3,m_4$，则这两组质量参数之间的关系为：
\begin{equation}
\begin{aligned}
m_a&=\frac12(m_1+m_2),\\
m_b&=\frac12(m_1-m_2),\\
m_c&=\frac12(m_3+m_4),\\
m_d&=\frac12(m_3-m_4).
\end{aligned}
\tag{21.9}
\end{equation}
对偶群 $\SL_2(\Z)$ 对参数 $m_{a,b,c,d}$ 的作用比对 $m_{1,2,3,4}$ 的作用更简单，因为它按式 (21.8) 所示的模式置换 $m_{a,b,c,d}$。例如，在 $S$ 下，$m_b$ 与 $m_c$ 交换，而 $m_a$ 与 $m_d$ 保持不变。可以
% [原文 p.338]
验证这一对偶作用对应于式 (20.14) 中的变换（在考虑式 (21.9) 的关系时）。

\noindent\textbf{边缘形变空间与对偶}

当所有超多重态都无质量时，库仑分支上的低能规范耦合为 $\tau=a_D(u)/a(u)=2\tilde\tau_{\mathrm{cl}}$，其中 $\tilde\tau_{\mathrm{cl}}$ 是重整化不变的 UV 复化规范耦合。用 UV 规范耦合 $g$ 与 theta 角 $\Theta$ 表示，有：
\begin{equation}
\tau=\frac{\Theta}{\pi}+\frac{8\pi\ii}{g^2}.
\tag{21.10}
\end{equation}
因此 $\tau$ 取值于复上半平面 $\HH$。对偶群 $\SL_2(\Z)$ 通过分式线性变换作用在 $\tau$ 上，诱导出 $\PSL_2(\Z)=\SL_2(\Z)/\Z_2$ 在 $\HH$ 上的标准作用。换言之，理论的规范耦合空间是 $\HH$，而模去 $\SL_2(\Z)$ 对偶的规范耦合空间是商 $\HH/\PSL_2(\Z)$。注意商 $\HH/\PSL_2(\Z)$ 只有一个尖点，对应于弱耦合极限 $\operatorname{Im}\tau\to\infty$。

\subsection{Seiberg–Witten 曲线}

带四个（有质量）基本超多重态的四维 $\mathcal{N}=2$ SU(2) 超 Yang–Mills（SYM）理论的 Seiberg–Witten 曲线由第 21.3 节得到：
\[
(\Sigma):\quad f\frac{(v-m_1)(v-m_2)}{t}+f'(v-m_3)(v-m_4)t=v^2-u.
\]
这里 $f$ 与 $f'$ 是复常数，$m_{1,2,3,4}$ 编码超多重态的裸质量，$u$ 是库仑分支上的坐标。Seiberg–Witten 微分由 $v\,\mathrm{d}t/t$ 给出。我们将改写 $\Sigma$，以突出 $\tau$ 的几何本质，这对讨论推广会有用。我们遵循 [339, Chap. 9] 的表述，更详细的细节可在那里找到。

可以重标度 $t$ 使 $f'=1$，并把 $v$ 的各次幂归并，得到：
\[
(\Sigma):\quad v^2+\frac{f(m_1+m_2)/t+(m_3+m_4)t}{1-t-f/t}v-\frac{u+fm_1m_2/t+m_3m_4t}{1-t-f/t}=0.
\]
令：
\[
x=v-\frac12\frac{f(m_1+m_2)/t+(m_3+m_4)t}{1-t-f/t},
\]
便得到：

% [原文 p.339]
\[
(\Sigma):\quad x^2-\left(\frac{u+fm_1m_2/t+m_3m_4t}{1-t-f/t}+\frac14\left(\frac{f(m_1+m_2)/t+(m_3+m_4)t}{1-t-f/t}\right)^2\right)=0.
\]
微分 $\lambda=x\,\mathrm{d}t/t$ 与 $v\,\mathrm{d}t/t$ 之间的关系只涉及质量参数，因此 $\lambda=x\,\mathrm{d}t/t$ 是 Seiberg–Witten 微分的另一个合适选择。

设 $t_1$ 与 $t_2$ 是 $1-t-f/t=-(t-t_1)(t-t_2)/t$ 的根，并令 $z=t/t_1$、$q=t_2/t_1$。则：
\begin{equation}
(\Sigma):\quad\lambda^2-\phi_2(z)=0,
\tag{21.11}
\end{equation}
其中
\[
\phi_2(z)=\frac{P(z)}{(z-1)^2(z-q)^2}\frac{\mathrm{d}z^2}{z^2}
\]
是 $\C\PP^1_z$ 上的一个二次微分，在 $z=0,q,1,\infty$ 处有二阶极点，而 $P(z)$ 是次数为 4 的多项式。$\PSL_2(\C)$ 在 $\C\PP^1_z$ 上的作用允许我们对称地处理二次微分 $\phi_2(z)$ 的四个奇点。

回忆 BPS 粒子的质量由 $|\lambda|$ 沿 $\Sigma$ 上闭链的积分计算。因此，$\lambda$ 的质量维数为 1，而 $P(z)$ 的质量维数为 2。超多重态的裸质量是 $\lambda$ 在 $\Sigma$ 上极点处的留数；可以求得 $\lambda$ 在 $z=0,q,1,\infty$ 上方各极点处的留数分别为 $\pm m_b,\pm m_a,\pm m_c,\pm m_d$\footnote{这在一个有限重整化之下成立，例如 [339] 中所讨论的。}。相应地，$\phi_2$ 在 $0,q,1,\infty$ 处二阶极点的系数分别为 $m_b^2,m_a^2,m_c^2,m_d^2$。例如：
\begin{equation}
\phi_2(z)=\frac{m_b^2}{z^2}\mathrm{d}z^2+\text{less singular terms},
\tag{21.12}
\end{equation}
在 $z$ 接近 0 时成立。因此，味道因子 $\SU(2)_a$、$\SU(2)_b$、$\SU(2)_c$、$\SU(2)_d$ 与 $\C\PP^1_z$ 上的穿孔 $0,q,1,\infty$ 之间存在自然的一一对应。注意 $\phi_2(z)$ 在 $\C\PP^1$ 上除这些穿孔外全纯。

% [Tutor Note: 质量维数核对：BPS 质量由 $\int|\lambda|$ 给出，故 $\lambda$ 的质量维数为 1；于是 $\phi_2=\lambda^2$ 的质量维数为 2。而 $\mathrm{d}z^2/z^2$ 无量纲，故 (21.12) 的极点系数 $m_b^2$ 质量维数恰为 2 ✓。这也解释了为什么极点的“系数”是质量的平方：$\lambda$ 的留数是 $\pm m_b$（质量维数 1），平方后即得 $\phi_2$ 二阶极点的系数。相应地 $\lambda$ 在四个穿孔上方的留数依次为 $\pm m_b,\pm m_a,\pm m_c,\pm m_d$，与图 21.9 的标签 $0\leftrightarrow b$、$q\leftrightarrow a$、$1\leftrightarrow c$、$\infty\leftrightarrow d$ 一致。]

\begin{remark}[21.16] 当所有超多重态都无质量（即 $m_{1,2,3,4}=0$）时，同样的计算给出形如式 (21.11) 的 Seiberg–Witten 曲线，其中
\[
\phi_2(z)=\frac{P}{(z-1)(z-q)}\frac{\mathrm{d}z^2}{z},
\]
而 $P$ 为常数。注意此时 $\phi_2(z)$ 在 $0,q,1,\infty$ 处有一阶极点。\end{remark}

\begin{definition}[21.17] 在无质量形变时 $\phi_2$ 为 1 阶、而在一般质量形变时为 2 阶的穿孔，称为简单 $\SU(2)$ 穿孔。\end{definition}

\begin{definition}[21.18] 带四个简单穿孔的黎曼球面是带四个基本味道的 $\SU(2)$ SYM 的 UV 曲线；它如图 21.9 所示。\end{definition}

\begin{remark}[21.19] 复数 $q\in\C\PP^1$ 是穿孔 $0,q,1,\infty\in\C\PP^1_z$ 的交比。交比在 Möbius 群 $\PSL_2(\C)$（黎曼球面的复自同构群）的作用下不变。\end{remark}

% [原文 p.340]
\Fig[0.5\linewidth]{figures/fig_21_9.png}{21.9}{带四个基本味道的 $\SU(2)$ SYM 的 UV 曲线。}

\begin{remark}[21.20] 上面 $q$ 的定义表明，$q$ 与常数 $f$ 有关，而 $f$ 本身与 UV 规范耦合 $\tilde\tau$ 有关，后者是理论的一个恰边缘参数——实际上是理论唯一的恰边缘参数。\end{remark}

UV 规范耦合 $\tilde\tau$ 自然地取值于复上半平面 $\HH$。事实上，$\HH$ 可等同于带四个穿孔的黎曼球面的泰希米勒空间（参看第 6 章）。

\begin{proposition}[21.21] 带四个基本味道的 $\SU(2)$ SYM 的 UV 曲线的泰希米勒空间，是理论的恰边缘参数空间。作用在 $\tilde\tau$ 上的对偶群，即 $\PSL_2(\Z)$，可以等同于 UV 曲线映射类群中非平凡地作用在泰希米勒空间 $\HH$ 上的子群\footnote{带四个穿孔的球面 $S_{0,4}$ 的映射类群是半直积 $\PSL_2(\Z)\ltimes(\Z_2\times\Z_2)$ [125]。}。商 $\HH/\PSL_2(\Z)$ 是带四个穿孔的球面 $S_{0,4}$ 上复结构的模空间。\end{proposition}

\begin{remark}[21.22] 这里，关于亏格为 $g$、带 $n$ 个穿孔的光滑曲面 $S_{g,n}$ 的泰希米勒空间，我们指的是 $S_{g,n}$ 所能装备的所有双曲度量的空间（其中穿孔是双曲尖点），模去同痕于恒等映射的微分同胚。\end{remark}

与无质量情形相反，当超多重态有质量时，$\HH$ 中的点与低能规范耦合之间的对应变得非平凡。尽管如此，式 (21.10) 中的 UV 规范耦合仍是理论的一个恰边缘参数。因此一般而言，$\HH=\mathcal{T}(S_{0,4})$ 应理解为理论的恰边缘形变空间。

纯映射类群是映射类群 $\operatorname{Mod}(S_{0,4})$ 中保持穿孔的特定子群（参看定义 4.17 与 [125]）。换言之，非平凡地作用在 $\HH$ 上的纯映射类群，是 $\PSL_2(\Z)$ 中保持质量 $m_{a,b,c,d}$ 的子群，因此可刻画为模 2 约化 $\PSL_2(\Z)\to\PSL_2(\Z_2)$ 的核 $\Gamma(2)$。所以，模去保持质量 $m_{a,b,c,d}$ 的对偶，理论的恰边缘形变空间是商 $\HH/\Gamma(2)$。

\begin{remark}[21.23] 商 $\HH/\Gamma(2)$ 是带四个标记穿孔的球面的模空间。\end{remark}

% [原文 p.341]
模空间 $\HH/\Gamma(2)$ 有三个尖点。它们与图 21.8 所示的箭图一一对应，从而与带四个基本味道的 $4d$ $\mathcal{N}=2$ SU(2) SYM（在质量参数有标记时）的三种不同弱耦合描述一一对应。几何上，模空间 $\HH/\Gamma(2)$ 中的三个尖点对应于四穿孔球面 $\Sigma_{0,4}$ 退化为两个球面（各含两个简单穿孔、由一根细长管子相连）的不同方式，如图 21.10 所示。关于这些弱耦合极限的显式描述，我们参见 [339, Chap. 9]。

\Fig[0.9\linewidth]{figures/fig_21_10.png}{21.10}{带四个穿孔的球面可能的退化。}

在每个这样的极限下，带四个穿孔的黎曼球面完全退耦为两个球面，每个都有三个简单穿孔。例如，极限 $\tau=\infty$ 对应于原 $\SU(2)$ 规范群的完全退耦，得到两组四个超多重态。因此，带三个穿孔的黎曼球面被认定为只含这四个超多重态的理论的 UV 曲线。

\subsection{Hitchin 描述}

式 (21.11) 可以改写为一种使与 Hitchin 系统的关系变得非常显明的形式。对应于二次微分 $\phi_2(s)$ 的 Hitchin 场是一个无迹 $2\times2$ 矩阵 $\Phi(z)$，满足 $\tr\Phi(z)^2=2\phi_2(z)$。用 $\Phi(z)$ 表示，式 (21.11) 改写为：
\begin{equation}
\det(\lambda\otimes\mathbf{1}_2-\Phi(z))=0.
\tag{21.13}
\end{equation}
注意式 (21.13) 在按共轭作用在 $\Phi(z)$ 上的 $\SU(2)$“规范”变换下不变，即 $\Phi(z)$ 在 $\SU(2)$ 的伴随表示中变换。在质量参数 $m$ 不为零的简单穿孔 $z=z_*$ 处，二次微分 $\phi_2(z)$ 有二阶极点，系数为 $m^2$（参看式 (21.12)）。等价地，在 $\SU(2)$ 规范变换之下，当 $z$ 接近 $z_*$ 时有：
\[
\Phi(z)\sim\begin{pmatrix}m&0\\0&-m\end{pmatrix}\frac{\mathrm{d}(z-z_*)}{z-z_*}.
\]

% [Tutor Note: 三点核对。(a) 原文“对应于二次微分 $\phi_2(s)$ 的 Hitchin 场”把自变量写成 $s$，但本段其余各处（以及 (21.11)–(21.13)）均用 $z$，按语境应为 $\phi_2(z)$，这是一处笔误，照译并在此指出。(b) 由 $\tr\Phi^2=2\phi_2$ 可推出 (21.13)：无迹 $2\times2$ 矩阵 $\Phi$ 的特征值为 $\pm a$ 时 $\tr\Phi^2=2a^2$、$\det(\lambda\mathbf{1}_2-\Phi)=\lambda^2-a^2=\lambda^2-\frac12\tr\Phi^2$，故 $\frac12\tr\Phi^2=\phi_2$ 恰给出 $\lambda^2-\phi_2=0$。(c) 极点结构：$\Phi\sim\operatorname{diag}(m,-m)\,\mathrm{d}(z-z_*)/(z-z_*)$ 时 $\tr\Phi^2=2m^2$，即 $\phi_2$ 在 $z_*$ 有二阶极点、系数 $m^2$，与 (21.12) 一致；当 $m\to0$ 时 $\Phi\sim\left(\begin{smallmatrix}0&1\\0&0\end{smallmatrix}\right)\mathrm{d}(z-z_*)/(z-z_*)$ 满足 $\tr\Phi^2=0$，故 $\phi_2$ 只剩一阶极点，正对应定义 21.17 的简单穿孔。]
当 $\lambda$ 的留数 $m$ 为零时，二次微分 $\phi_2$ 在 $z_*$ 处极点的阶降为一。在规范变换之下，当 $z$ 接近 $z_*$ 时可以假设

% [原文 p.342]
\[
\Phi(z)\sim\begin{pmatrix}0&1\\0&0\end{pmatrix}\frac{\mathrm{d}(z-z_*)}{z-z_*}.
\]
注意 Hitchin 场 $\Phi(z)$ 在简单穿孔处有一阶极点，无论是否考虑质量形变。关于 S 类理论的 Hitchin 系统与 Seiberg–Witten 纤维化之间关系的更多细节，读者例如可参阅 [155]。

\subsection{带 $F<4$ 个味道的 SU(2) SYM}

把上述推理应用于带 $F<4$ 个味道的 SU(2) SYM 的 Seiberg–Witten 曲线是富有启发性的。对 $F=0$，Seiberg–Witten 曲线读作
\[
\frac{\Lambda^2}{t}+\Lambda^2t=v^2-u.
\]
可把它改写为：
\begin{equation}
\lambda^2-\phi_2(z)=0,
\tag{21.14}
\end{equation}
其中 $\lambda=x\,\mathrm{d}z/z$，$\phi_2$ 是球面 $\C\PP^1_z$ 上的二次微分，在零与无穷处有三阶极点。转而考虑 $F=1$ 的情形，它由最左边 NS5-膜左侧带一张半无限 D4-膜的 IIA 型构造得到，此时 Seiberg–Witten 曲线读作：
\[
\frac{2\Lambda_1(v-m)}{t}+\Lambda_1^2t=v^2-u.
\]
可把它改写为式 (21.14) 的形式，其中 $\phi_2$ 在 0 处有 4 阶极点、在 $\infty$ 处有 3 阶极点。还有一种把 0 与 $\infty$ 处穿孔的角色互换的类似描述；它由最右边 NS5-膜右侧带半无限 D4-膜的 IIA 型构造开始得到。

$F=2$ 与 $F=3$ 的情形也允许类似的描述。例如，当 $F=2$ 时，微分 $\phi_2$ 可以取为在 0 与 $\infty$ 各有一个 4 阶极点，或者在 0、1 处各有一个简单极点、在 $\infty$ 处有一个 3 阶极点。即使在带四个基本味道的 SU(2) SYM 情形，也可以考虑左右两侧半无限 D4-膜数目不匹配的非平衡 IIA 型描述。微分 $\phi_2(z)$ 的极点阶数大于 2 的穿孔称为不规则 $\SU(2)$ 穿孔。$\phi_2$ 的更高阶极点编码 NS5-膜的弯曲，后者由终止于其两侧的 D4-膜数目不平衡所引起。

在 Seiberg–Witten 纤维化的 Hitchin 描述中，$\phi_2$ 有高于二阶极点的这类穿孔对应于 Hitchin 场的不规则高阶极点。这引出下面的区分，它将推广到其他规范群：
